\section{A compositional finite-observation framework}
\label{sec:compiler}

\subsection{Finite safety under infinite hypotheses}

A frequent mistake is to justify an induction by saying that every
individual computation is finite. The induction formula itself must be
permitted. The following lemma gives the needed shape.

\begin{theorem}[Finite safety compilation]
\label{thm:safety}
Fix primitive recursive functions $\operatorname{step}(\mathbf x,j,s)$,
$s_0(\mathbf x)$, a horizon $H(\mathbf x)$, and a Boolean invariant
$I(\mathbf x,j,s)$. Let $\Gamma(\mathbf x)$ be any fixed collection of
first-order assumptions. Suppose there are \iPRA{} derivations, under
$\Gamma$, of
\[
 I(\mathbf x,0,s_0)=1,
 \qquad
 j<H\ \wedge\ I(\mathbf x,j,s)=1
 \ \longrightarrow\ I(\mathbf x,j+1,\operatorname{step}(\mathbf x,j,s))=1.
\]
Define $s_{j+1}=\operatorname{step}(\mathbf x,j,s_j)$ by primitive
recursion. Then there is an \iPRA{} derivation under $\Gamma$ of
\[
 \forall j\le H\quad I(\mathbf x,j,s_j)=1.
\]
This is a proof-transformation theorem: the base and step derivations
are required, not inferred merely from their truth in $\N$.
\end{theorem}
\begin{proof}
Keep $\mathbf x,H$ fixed. Induct on the quantifier-free formula
\[
 j\le H\ \longrightarrow\ I(\mathbf x,j,s_j)=1.
\]
At a successor, a rational/numerical case split first checks $j+1\le H$.
If it fails the implication follows. If it holds, $j<H$ and $j\le H$;
the induction hypothesis and the supplied step derivation give the
conclusion. The hypotheses $\Gamma$ remain in the context. Quantifier-free
induction does not require the other context formulas to be quantifier-free.
Ordinary universal introduction subsequently closes the numerical
parameters. Syntactically, encode this guarded formula as a
\texttt{Fm PTm}; its \texttt{QF} proof permits the single application
of \texttt{PRA.qfInd}.
\end{proof}

A finite contract can accumulate many facts: bounds on a state,
well-formedness of a list, every earlier local transition check,
a polynomial coefficient identity, and every relevant approximation
inequality. A bounded conjunction of those facts becomes one Boolean
invariant by Lemma~\ref{lem:reflection}. A universal name hypothesis is
used only to verify the finitely many input observations required by that
particular run.

\subsection{Two different budgets}

An observation procedure has two potentially different resource bounds:
\begin{enumerate}
\item its \emph{evaluation demand} is the largest input observation needed
      to compute an output observation;
\item its \emph{verification demand} is the largest input observation and
      contract instance needed to prove the desired finite output facts.
\end{enumerate}
The second can be larger. A root approximation at precision $k$ might
query one finite prefix, whereas comparing it to another approximation or
proving a residual inequality asks for additional rows. An implementation
trace alone does not bound its mathematical proof obligations. We carry
both bounds and use their maximum when composing certified operations.

For simplicity write all observations as one sequence indexed by $j$.
A tuple index $(i,k)$ can be flattened using a fixed numerical coding, or
one can keep rectangular tables with the same argument. Let
$\operatorname{Tab}(a,h)=[a(0),\ldots,a(h)]$. The zero extension of a
finite table gives a total input to any numerical program; correctness
conditions guarantee that no relevant access uses the extension.

\begin{definition}[Finite local contract]
\label{def:contract}
A stage $s$ of a construction is specified by:
\begin{itemize}
\item a primitive recursive numerical table program $F(s,p,X)$;
\item an evaluation bound $u(s,p)$ such that observation $p$ uses only the
      input entries indexed at most $u(s,p)$;
\item a Boolean finite property $C_s(h,X)$, using only the prefix through
      $h$ and downward closed under restriction to shorter prefixes;
\item a primitive recursive verification bound $v(s,h)$ and a proved
      finite implication
\begin{equation}
 C_s(v(s,h),X)=1
 \quad\Longrightarrow\quad
 C_{s+1}\bigl(h,T_s^h(X)\bigr)=1.             \label{eq:contract}
\end{equation}
\end{itemize}
Here $T_s^h(X)$ is the table of outputs $0\le p\le h$, evaluated from the
appropriate prefixes of $X$. The bound $v$ dominates the evaluation
bounds for those outputs. The predicates $C_s$ are values of one fixed
primitive recursive predicate $C(s,h,X)$. There is one \iPRA{}
derivation of \eqref{eq:contract} with $s,h$ and the metadata as free
numerical parameters, ordinarily obtained by Theorem~\ref{thm:safety}.
A separate external proof for every fixed $s$ does not meet this
uniformity requirement.
\end{definition}

Metadata can vary with the stage. It is included among the numerical
parameters and either computed exactly in advance or bounded by a fixed
primitive recursive recurrence. Metadata affecting the outputs is fixed
canonically from the original inputs and stage, independently of the
requested horizon, or is covered by the same prefix-locality condition.
In the bounded-metadata case replace the demands
by their bounded maxima over all metadata codes below the supplied cap;
substituting a cap into a nonmonotone demand is not sufficient.
If a particular interface requires extra output facts,
include their finite tests in $C_s$. We do not claim that every mathematical
property automatically has such a finite contract.

\begin{theorem}[Compilation of variable-depth observation constructions]
\label{thm:demand}
Suppose a fixed family of stages has the data of
Definition~\ref{def:contract}, uniformly primitive recursive in its
numerical stage and metadata. For any numerical number $N$ of stages
and desired final observation horizon $k$, there is one primitive
recursive table computation producing all final observations through
$k$ from a bounded prefix of the initial presentation. If that prefix
satisfies its initial finite contract, the output satisfies its final
finite contract. In applications a uniform derivation from the original
assumptions verifies the initial contract. The statement is uniform
in $N$ and $k$.

In particular, repeated application of \emph{one fixed} certified
operation does not require primitive recursion at function type, or a
universal evaluator for all primitive recursive programs.
\end{theorem}
\begin{proof}
With the starting metadata and $N$ fixed as numerical parameters, define
\[
 U_N(h)=\max\bigl( h,\{u(s,p):s<N,\ p\le h\},
                         \{v(s,p):s<N,\ p\le h\}\bigr).
\]
For $N=0$ the empty maxima contribute zero. This is a primitive recursive
function, nondecreasing in $h$, with $U_N(h)\ge h$. Stage-dependent metadata
is first computed by its fixed numerical bound recurrence and included in
the corresponding values of $u,v$.

Compute $d_0=k$ and $d_{r+1}=U_N(d_r)$ for $r<N$. For stage $s$ put
\begin{equation}
 h_s=d_{N-s}\qquad(0\le s\le N).               \label{eq:budgets}
\end{equation}
Thus $h_N=k$ and $h_s=U_N(h_{s+1})$. Initially tabulate the fixed input
observation function through $h_0$. Inductively form a new table
\begin{equation}
 X_{s+1}[p]=F\bigl(s,p,
          \operatorname{restrict}(X_s,u(s,p))\bigr),
 \qquad 0\le p\le h_{s+1}.                   \label{eq:table-build}
\end{equation}
The complete table is built by a bounded loop; the outer recursion is on
$s$ with $N,k$ and the already computed list of budgets fixed. Its state
is a \emph{number coding a finite table}. Every requested entry exists
since $u(s,p)\le h_s$.

All these operations are first-order primitive recursive. Nested
precision requests have become a forward computation of numerical
budgets followed by a forward computation of numerical tables. Even if
$U_N$ produces enormous exponent towers, its $N$-fold numerical iterate
is a primitive recursive function of $N,k$ and the starting metadata.
The number of stored observations is at most
$(N+1)(d_N+1)$. We do not claim a useful polynomial running time.

For the contract proof, use as the bounded invariant
\[
 s\le N\ \longrightarrow\ C_s(h_s,X_s)=1.
\]
The initial contract is a finite consequence of the input hypotheses.
Restriction supplies $C_s(v(s,h_{s+1}),\cdot)=1$ and
\eqref{eq:contract} gives the next row. Both the row construction and the
invariant are numerical. Theorem~\ref{thm:safety} applies, so there is
no induction on an unbounded statement about an entire output name.
\end{proof}

\paragraph{Coherence across requested horizons.}
Computing with two different final horizons must not change an output
already computed. This follows from deterministic prefix locality: if
$X,X'$ agree on every entry through $u(s,p)$, then their computed output
at $p$ is identical. An induction on the finite stage count shows
agreement of two table runs on their common relevant rows. Consequently
\[
 A_N(p):=X_N^{(N,p)}[p]
\]
is a single primitive recursive observation function. When a proof needs
several of its values simultaneously, compute one table at the maximum
verification horizon; locality identifies its entries with the individual
definition. This last step is essential: unrelated finite output tables
would not define a single observation procedure. Prefix consistency by
itself is not Cauchy regularity. When a regular name is required, its
pairwise finite regularity inequalities must be included in $C_N$;
arbitrary-horizon introduction then yields the full regularity schema.

The checked Lean theorem \texttt{compiled\_correct} verifies precisely
this comparison between the finite table implementation and the nested
semantic construction. Its companion \texttt{ofBody\_local} derives
locality from a numerical prefix-table body. The semantic construction
uses function-valued recursion only as an ambient Lean specification;
\eqref{eq:table-build} is the first-order arithmetic implementation.

\begin{corollary}[Eliminating function-valued iteration in a certified fragment]
\label{cor:first-order}
A construction written informally as
$A_{s+1}=T_s(A_s)$ admits a uniform first-order primitive recursive
implementation whenever its stage programs, metadata bounds, evaluation
demands, and finite verification contracts satisfy the theorem. Exact
algebraic operations, finite probability transformations, formal-series
coefficient operations, and precision-indexed analytic constructions can
share this same compiler.
\end{corollary}

\subsection{Finite error transport through a chain}

Many operations are verified only up to an arbitrarily small rational
error. A single error estimate can prevent induction on exact equality
of represented objects.

\begin{lemma}[Telescoping defects]
\label{lem:defects}
In a normed ring of finite rational data with a submultiplicative norm,
let $P_0,\ldots,P_N$ and $A_0,\ldots,A_{N-1}$ satisfy
\[
 \|P_s-A_sP_{s+1}\|\le\varepsilon_s,
 \qquad \|A_s\|\le G_s.
\]
Then
\begin{equation}
 \left\|P_0-\left(\prod_{s<N}A_s\right)P_N\right\|
 \le \sum_{s<N}\varepsilon_s\prod_{j<s}G_j.     \label{eq:defectsum}
\end{equation}
If every $\varepsilon_s\le\varepsilon$ and $G_s\le G$ with $G\ge1$,
the right side is at most $NG^N\varepsilon$.
\end{lemma}
\begin{proof}
For each $r\le N$, let $B_r=A_0\cdots A_{r-1}$, with $B_0=1$.
Subtract consecutive terms in $B_rP_r$:
\[
 B_rP_r-B_{r+1}P_{r+1}=B_r(P_r-A_rP_{r+1}).
\]
Summing this exact finite identity gives the difference on the left of
\eqref{eq:defectsum}. Repeated triangle inequality and
submultiplicativity give the asserted bound. The proof is induction on
finite coefficient lists, finite products, and finite sums; at no point
is the induction formula an equality of infinite names. For polynomials
use the coefficient $\ell^1$ norm, whose submultiplicativity is just the
finite convolution triangle inequality. If complex magnitudes are used
in explanatory notation, rational upper bounds or squared-norm
certificates give the finite numerical formulation.
\end{proof}

To prove an exact represented identity, first fix an output precision
$p$, then select a local tolerance satisfying
$NG^N\varepsilon\le2^{-p}$ and compile only the observations required at
that tolerance. Universal introduction in $p$, or
Lemma~\ref{lem:vanishing-error}, supplies the exact relation. This order
of operations also works for matrix products, compositions of approximate
isometries, telescoping integrals, and finite sums of probabilistic errors.

\subsection{Why the conditions are a real restriction}

There cannot be a total primitive recursive function $E(e,n)$ which
evaluates every unary primitive recursive function coded by $e$: if there
were, $d(n)=E(n,n)+1$ would itself be primitive recursive, with a code
$e_d$, and $d(e_d)=E(e_d,e_d)+1=d(e_d)+1$ would be impossible.
This is an external diagonal argument, not an assertion that \iPRA{}
proves its own soundness.

A uniformly bounded machine-step interpreter is different. Given a
program, its input, and an explicit fuel bound, one can perform that many
primitive recursive steps and return either an answer or a timeout. For
an individual fixed primitive recursive program, its numerical resource
bounds can be proved by its particular finite construction. Neither fact
creates a primitive recursive universal total evaluator.

Likewise, a merely pointwise statement that every stage eventually stops
does not satisfy Theorem~\ref{thm:demand}. The crucial input is the
\emph{uniform primitive recursive} numerical demand recurrence and its
proof. There is also no automatic conversion of an arbitrary Lean
correctness proof into a \texttt{PRAAx} proof. The finite safety theorem
requires actual base and step derivations in the object theory.

\subsection{Relation to proof interpretation}

The framework follows the established idea of extracting quantitative
content and majorants, but specializes it to explicit first-order
numerical compilation. Functional interpretation and majorization have a
much broader literature \cite{kohlenbach-notes,kohlenbach-point}. Their
higher-type settings must not be imported silently. In particular,
Kohlenbach's analysis of uniform weak K\"onig's lemma shows that apparently
small changes involving uniform selection and extensionality can greatly
increase proof-theoretic strength \cite{kohlenbach-uwkl}. Our applications
below exhibit the finite contracts directly and do not appeal to an
unspecified conservation theorem.
